% §18.4 Estimation under Adversarial Queries
%
% Monograph register: the estimation problem, why iid fails, MAB framing
% (UCB, Thompson), CI under adversarial sampling (union bound + the
% towell-es-conf-mab bootstrap), connection to ch11/ch16, worked example.
% Tightened.
% FIX: the union-bound widening was stated as "tightened by roughly
% log(2e6) ~ 14.5 standard deviations", conflating ln(K/delta) with a
% standard-deviation count. The Bonferroni correction takes per-arm
% confidence to 1 - delta/K, so the Gaussian quantile rises from ~1.96 to
% sqrt(2 ln(K/delta)) ~ 5.4 at K=1e5, delta=0.05, widening each interval
% by about 2.7x. Corrected.
% Cross-refs sec:binomial-estimator, sec:confidence-intervals, thm:chernoff,
% sec:sample-complexity, sec:adversarial-inputs, sec:confidentiality-measures,
% sec:toy-system. Cites towell-es-conf-mab. DEFINED: sec:adversarial-estimation.

\section{Estimation under Adversarial Queries}
\label{sec:adversarial-estimation}

\Cref{sec:confidentiality-measures} priced confidentiality from $\fprate$,
$\fnrate$, and the workload prevalence $\pi$. The first two are fixed at
deployment; the third is a property of how the system is used, and a
defender who wants the actual leakage must estimate $\pi$ from live
traffic. The problem looks elementary, each query matching or not, the
matches Bernoulli, the chapter 7 machinery pricing a confidence interval.
The complication is the source of the queries. A benign iid workload is
the textbook setting of \Cref{sec:binomial-estimator}, and the Wilson and
Clopper-Pearson intervals of \Cref{sec:confidence-intervals} deliver a CI
in one line. An adversarial workload, where the attacker chooses each
query to maximize leakage and adapts to responses, is neither iid nor
stationary: the per-query match probability is a function of the
attacker's strategy, not a fixed parameter. The iid CI then bounds the
empirical match rate the defender just measured, not the worst-case rate
the attacker can sustain, the wrong quantity to feed the confidentiality
formula.

\subsection*{Multi-armed bandit framing}

The right framing treats each distinct keyword the attacker might query as
an arm in a multi-armed bandit\index{multi-armed bandit}: each pull is one query, the reward is the
leakage it produces, and the defender's goal is not to maximize reward but
to bound the worst arm's, a high-confidence statement that no keyword
exceeds a target leakage budget. The upper-confidence-bound family maps
directly: the defender maintains a per-arm estimate $\hat\pi_w$ and the UCB
$\hat\pi_w + \sqrt{2\log(1/\delta)/n_w}$, a high-confidence upper bound
that remains valid when the attacker chooses which arm to pull, because
it conditions on each arm's own samples rather than on an iid selection
rule. (With adaptively chosen sample sizes the fixed-$n$ radius needs a
union bound over sample sizes, or an anytime-valid radius, to keep its
coverage; the per-arm samples must also be fresh draws, not repetitions
of one frozen answer.) The
maximum per-arm UCB is the worst-case bound fed into the confidentiality
formula. Thompson sampling is the Bayesian counterpart, a
$\mathrm{Beta}(\alpha_w, \beta_w)$ posterior per arm with the maximum
sampled rate reported; both agree that the worst arm is what must be
bounded.

\subsection*{Confidence intervals under adversarial sampling}

The Wilson and Clopper-Pearson CIs\index{confidence interval!Wilson}\index{confidence interval!Clopper--Pearson} adapt by a union bound over arms. For a
vocabulary of $K$ arms at joint confidence $1 - \delta$, computing each
per-arm CI at $1 - \delta/K$ and taking the maximum is valid by
Bonferroni. The price is a $\log K$ factor in the per-arm width: the
Gaussian quantile rises from about $1.96$ at the marginal $\delta = 0.05$
to $\sqrt{2 \ln(K/\delta)}$ under the correction, which at $K = 10^5$ and
$\delta = 0.05$ is $\sqrt{2 \ln(2 \times 10^6)} \approx \sqrt{29} \approx
5.4$, widening each interval by roughly a factor of $2.7$. A tighter
alternative is the construction \textcite{towell-es-conf-mab} developed
for this setting: the defender estimates not $K$ independent quantities
but a single one, the worst-case confidentiality as the maximum of $K$
rates, and applies an order-statistic-aware moving-average bootstrap,
resampling the empirical trace with a moving window and reporting the
worst-window rate with the bootstrap quantile as the tail probability. The
bootstrap CI is narrower than the union bound when the per-arm rates are
similar and degenerates to it when one arm dominates. The defender chooses
the construction from what is known of the attacker: the union bound under
a worst-case assumption, the bootstrap when the attacker concentrates on a
few keywords.

\subsection*{Why the iid bound is the wrong bound}

The chapter 11 multiplicative Chernoff bound\index{Chernoff bound} (\Cref{thm:chernoff})
sharpens the iid CI to exponential decay, the right bound for iid traffic
and the wrong one for adversarial traffic, because its derivation invokes
the iid moment generating function, which assumes the arm-selection rule
is state-independent. The UCB bound replaces that MGF with an adversarial
concentration argument surviving state-dependent selection; the two agree
on iid traffic (the UCB constant $\sqrt{2\log(1/\delta)}$ matches the
Chernoff exponent for moderate $\delta$) and diverge on adversarial
traffic, where Chernoff under-bounds the true rate and UCB over-bounds it
correctly. The chapter 16 boundary discussion (\Cref{sec:adversarial-inputs})
reads here as a constructive recipe: the Bernoulli model itself is correct
(the per-query match indicator is still Bernoulli given the queried
keyword); what changes is the estimator aggregating the indicators into a
confidentiality bound.

\subsection*{A worked example}

A defender monitors a system at 100 queries per minute, processing $n =
6000$ queries in an hour and observing $540$ matches, an empirical rate
$\hat\pi = 0.09$. The \emph{iid} normal-approximation interval at $95\%$ is
\[
  0.09 \pm 1.96 \sqrt{0.09 \cdot 0.91 / 6000} \;\approx\; 0.09 \pm 0.0072,
\]
so $\pi \in [0.083, 0.097]$; substituting the upper end into the
confidentiality formula at $\fprate = 0.01$, $\fnrate = 0$ gives $\PPV
\approx 0.915$, $C \approx 0.085$, so the defender concludes
confidentiality is at least $0.085$. The \emph{MAB UCB} approach, with $K
= 1000$ keywords queried over the hour at $\delta_{\mathrm{arm}} =
0.05/1000 = 5 \times 10^{-5}$ and an average $n_w = 6$ queries per arm
(with high variance), bounds an arm at $n_w = 6$, $\hat\pi_w = 0.5$ by a
Clopper-Pearson upper limit near $1$ (about $0.99$ at this sample size and
joint confidence); the worst-arm bound is then $\pi^{\mathrm{MAB}}_{\max}
\approx 0.99$, giving $\PPV \approx 0.9999$ and $C \approx 10^{-4}$. The two cover estimates differ by nearly three orders of magnitude, and
only the MAB bound answers the defender's actual question, how leaky is the worst-case keyword the
attacker might target. The warning is the chapter's recurring one: a
framework guarantee is only as strong as the estimator that produced its
parameters, and the iid estimator misapplied to adversarial data, not the
Bernoulli machinery, is at fault. The next section turns from estimation
back to construction.
